\section{Introducción}

\subsection{Objetivos del trabajo práctico}
El presente trabajo práctico tuvo como objetivo determinar la conexión de un transformador para que funcione como elevador o como reductor de tensión, comprobar el efecto de aislamiento eléctrico del transformador y Analizazr el efecto que tiene una carga en un circuito con un transformador.

\subsection{Circuito analizado}
El esquemático correspondiente al circuito estudiado en la etapa principal de análisis se detalla en la Figura~\ref{fig:circuito_trafo}. Este consistió en un generador de señales de corriente alterna que provee una tensión senoidal $V_S$ con una frecuencia de $f = 50\text{ kHz}$. En serie con la fuente se conectó una resistencia de medición $R_m = 47\ \Omega$, acoplada al devanado primario del transformador (de inductancia $L_1$). El devanado secundario (de inductancia $L_2$) cerró la malla con una resistencia de carga de valor $R_L = 1\text{ k}\Omega$. Las mediciones se realizaron alternando el sentido de conexión del transformador para registrar las variaciones de la impedancia de entrada y las tensiones del sistema mediante el uso de instrumental de laboratorio.

\begin{figure}[htbp]
\centering
\begin{circuitikz}[american, scale=1.1]

    % Transformador
    \draw (2,3) node[transformer core, anchor=A1] (T) {};

    % Fuente Vs - lado izquierdo
    \draw (T.A1) to[R, l=$R_m$, -*] ++(-2,0)
                 to[sinusoidal voltage source, l=$V_S$] ++(0,-2.1)
                 -- (T.A2);

    % Etiquetas de inductancias
    \node[below] at (T.A2 -| T.west) {$L_1$};
    \node[below] at (T.B2 -| T.east) {$L_2$};

    % Carga RL - lado derecho
    \draw (T.B1) to[short, -*] ++(1,0) coordinate(topR)
                 to[R, l=$R_L$] ++(0,-2.1) coordinate(botR)
                 -- (T.B2);

    % Puntos de homología (lado interior, arriba de cada bobina)
    \node[circle, fill, inner sep=1.5pt] at ($(T.A1)!0.15!(T.A2)+(0.45,0)$) {};
    \node[circle, fill, inner sep=1.5pt] at ($(T.B1)!0.15!(T.B2)+(-0.45,0)$) {};

    % Tierra bajo Vs (nodo inferior izquierdo)
    \draw (T.A2) ++(-2,-0.0) to[short] ++(0,-0.3) node[ground]{};

    % Tierra bajo RL (nodo inferior derecho)
    \draw (botR) to[short] ++(0,-0.3) node[ground]{};

\end{circuitikz}
\caption{Circuito de análisis para el ensayo del transformador acoplado con carga.}
\label{fig:circuito_trafo} % <-- Colocado abajo del caption para fijar el número correlativo correcto
\end{figure}

\subsection{Marco teórico}
El principio básico de funcionamiento de un transformador se fundamenta en la ley de inducción de Faraday y en el fenómeno de inducción mutua entre dos circuitos vinculados por un flujo magnético común. Cuando una corriente variable en el tiempo circula por el devanado primario, genera un flujo magnético que concatena al devanado secundario, induciendo en este último una fuerza electromotriz (fem) proporcional a la tasa de cambio del flujo y al número de espiras.

En un \textbf{transformador ideal}, se asume que la permeabilidad magnética del núcleo es infinita ($\mu \rightarrow \infty$), las inductancias tienden a infinito, el acoplamiento es perfecto ($K = 1$) y no existen pérdidas óhmicas ni en el núcleo ($R_1 = R_2 = 0$). Bajo estas hipótesis, las tensiones y corrientes se rigen linealmente por la relación de transformación $n$:
\begin{equation}
    n = \frac{N_2}{N_1} = \frac{V_2}{V_1} = \frac{I_1}{I_2}
\end{equation}

Por el contrario, un \textbf{transformador no ideal o real} presenta limitaciones físicas que alteran este comportamiento. El acoplamiento entre bobinas nunca es perfecto ($K < 1$), lo que da lugar a flujos de dispersión modelados matemáticamente mediante inductancias de dispersión primarias ($L_{d1} = L_1(1-K)$) y secundarias referidas ($L_{d2}' = L_2'(1-K)$). Adicionalmente, se deben considerar las resistencias óhmicas de los conductores de cobre ($R_1$ y $R_2$), y las pérdidas en el núcleo magnético debidas a ciclos de histéresis y corrientes parásitas de Foucault. Estas se representan en el circuito equivalente por medio de una conductancia de pérdidas $G_0$ en paralelo con la inductancia de magnetización $L_0$.

Para el análisis cuantitativo del transformador real, el coeficiente de acoplamiento magnético ($K$) permite evaluar la fracción del flujo magnético total que concatena a ambos devanados. Este coeficiente adimensional se encuentra acotado en el rango de $0 \le K \le 1$ y se define matemáticamente a partir de la inductancia mutua ($M$) y las inductancias propias de cada puerto ($L_1$ y $L_2$) mediante la expresión:
\begin{equation}
    K = \frac{M}{\sqrt{L_1 \cdot L_2}}
\end{equation}

\begin{itemize}
    \item \textbf{Caso Ideal ($K = 1$):} Representa un acoplamiento perfecto donde no existen flujos de dispersión; todo el flujo generado por un devanado atraviesa íntegramente al otro. En esta condición ideal, y asumiendo un núcleo sin pérdidas, la relación de transformación matemática coincide exactamente con la raíz cuadrada de la relación de inductancias propias:
    \begin{equation}
        n = \sqrt{\frac{L_2}{L_1}}
    \end{equation}
    \item \textbf{Caso Real ($0 < K < 1$):} El acoplamiento es imperfecto y aparecen inductancias de dispersión. En este escenario, la relación anterior deja de cumplirse de manera exacta para las tensiones terminales debido a las caídas internas en las reactancias de dispersión y resistencias serie.
    \item \textbf{Caso Desacoplado ($K = 0$):} Las bobinas se encuentran lo suficientemente distanciadas o blindadas como para que no exista inducción mutua entre ellas ($M = 0$).
\end{itemize}

Por otro lado, para simplificar la resolución de circuitos acoplados mediante mallas y nodos, se recurre frecuentemente a la \textbf{proyección o reflexión de parámetros} de un puerto a otro. Al analizar el circuito equivalente referido al devanado primario (Puerto 1), cualquier impedancia de carga $Z_L$ conectada en el devanado secundario (Puerto 2) se proyecta hacia la entrada modificada por la relación de espiras. Las ecuaciones fundamentales de transferencia y proyección entre ambos puertos se definen como:
\begin{equation}
    V_2' = \frac{V_2}{n}
\end{equation}
\begin{equation}
    I_2' = n \cdot I_2
\end{equation}
\begin{equation}
    Z_L' = \frac{Z_L}{n^2}
\end{equation}

Estas relaciones de proyección permiten sustituir el acoplamiento inductivo por un circuito galvánicamente conectado equivalente, facilitando la determinación analítica de la impedancia de entrada total ($Z_e$) vista desde los terminales del generador de señales.

\newpage
